% Fragment: fig-sealed-two-adversaries.tex -- the same ledger mechanism,
% read against two different adversaries, paired accounting panels.
%
% Reader question: which of the two accounting mechanisms, the privacy
% ledger or the bit-capacity budget, actually bounds a given adversary?
% Claim: for an honest worker the ledger's sigma <= epsilon_max bound
% governs; for a malicious worker choosing its own reported spend, that
% same inequality certifies nothing, and the q*b capacity account of
% Sec. sealed-budget governs instead.
% Evidence: "Whom the certificate binds" paragraph, website-v2/public/
% whitepaper/sealed-harbor.tex, Sec. sealed-conservation ("Against the
% malicious worker...such a worker is not running a private mechanism...
% sigma <= epsilon_max certifies nothing, and the q*b capacity account...
% is the operative bound. The two accounts are for two different
% adversaries and the chapter owns both.").
% Must distinguish: the ledger mechanism itself is drawn IDENTICALLY in
% both panels (same guard, same arithmetic) -- what differs is only
% whether the recorded epsilon_i can be trusted to be what it claims;
% the figure must not suggest the ledger behaves differently for the two
% adversaries, only that its guarantee MEANS something different.
% Grammar chosen: paired accounting panels, left honest / right malicious,
% joined by one shared label naming both as one design -- the register
% row's own recommended form.
% Grammar rejected: a single flowchart with an if/else on worker honesty
% -- that draws a runtime branch the ledger itself never takes (the
% ledger runs the same code either way); the difference this figure must
% show is in what the SAME output means, not in what the mechanism does.
% Five-second test: the two ledger boxes are identical; only the verdict
% below each differs.
% QA: beauty_lint warns B7 in swiss/technical (sub-point near-alignment
% noise between the two panels' mirrored node centres); justified, not
% fixed.
\begin{figure}[H]
\centering
\pdfigurehue{pdcobalt}%
\begin{tikzpicture}[pd figure,x=1cm,y=1cm]
  \def\dy{1.55}
  \def\xA{0}   \def\xB{6.0}
  \node[pd focus state,minimum width=3.55cm,align=center] (WA) at (\xA,{2*\dy})
    {runs an $\varepsilon_i$-private mechanism};
  \node[pd breach state,minimum width=3.55cm,align=center] (WB) at (\xB,{2*\dy})
    {declares its own $\varepsilon_i$};
  \node[pd state,minimum width=3.55cm,align=center] (LA) at (\xA,{1*\dy})
    {ledger: $\sigma\le\varepsilon_{\max}$?};
  \node[pd state,minimum width=3.55cm,align=center] (LB) at (\xB,{1*\dy})
    {ledger: $\sigma\le\varepsilon_{\max}$?};
  \node[pd ready state,minimum width=3.55cm,align=center] (VA) at (\xA,0)
    {ledger governs: $(\varepsilon_{\max},0)$-DP};
  \node[pd breach state,minimum width=3.55cm,align=center] (VB) at (\xB,0)
    {certifies nothing};
  \node[pd ready state,minimum width=3.55cm,align=center] (QB) at (\xB,{-1*\dy})
    {$q\cdot b$ governs instead};
  \draw[pd focus arrow] (WA) -- (LA);
  \draw[pd breach arrow] (WB) -- (LB);
  \draw[pd focus arrow] (LA) -- (VA);
  \draw[pd breach arrow] (LB) -- (VB);
  \draw[pd focus arrow] (VB) -- (QB);
  \begin{scope}[on background layer]
    \node[pd panel,fit={(WA)(LA)(VA)},inner xsep=9pt,inner ysep=9pt] (PA) {};
    \node[pd panel,fit={(WB)(LB)(VB)(QB)},inner xsep=9pt,inner ysep=9pt] (PB) {};
  \end{scope}
  \node[pd title,anchor=south] at (PA.north) {honest worker};
  \node[pd title,anchor=south] at (PB.north) {malicious worker};
  \draw[pd guide] (PA.north) ++(0,0.62) -- ++(0,0.55) -| node[pd tag,pos=0.5,anchor=south]
    {one design, two adversaries} ($(PB.north)+(0,0.62)$);
\end{tikzpicture}
\caption{The same ledger code, read against two adversaries. Left, an
honest worker's recorded $\varepsilon_i$ is the true cost of an
$\varepsilon_i$-private mechanism, so the ledger's guard
$\sigma\le\varepsilon_{\max}$ licenses genuine $(\varepsilon_{\max},0)$-differential
privacy (Theorem~\ref{thm:sealed-conservation}). Right, a malicious worker
declares whatever $\varepsilon_i$ it likes, so the identical guard
certifies nothing about real leakage; the operative bound for that
adversary is instead the $q\cdot b$ capacity account of
Sec.~\ref{sec:sealed-budget}. The ledger box is drawn identically in both
panels because it runs identically in both cases --- only what its output
is entitled to mean changes. [internal]}
\label{fig:sealed-two-adversaries}
\end{figure}
